%%%%%%%%%%%%%%%%%%%%%%%%%%%%%%%%%%%%%%%%%%%%%%%%%%%%%%%%%%%%%%%%%%%%%%
% Lớp 4 | Chủ đề 1. Ôn tập và bổ sung | Số chẵn, số lẻ
% Dạng 4D1?3-2: Tính chẵn lẻ của tổng, hiệu, tích
%%%%%%%%%%%%%%%%%%%%%%%%%%%%%%%%%%%%%%%%%%%%%%%%%%%%%%%%%%%%%%%%%%%%%%


% nguồn: 2026-11-18_lop4_chu-de-1-tu-luan (nhập 18/11/2026)
\begin{bt}%[4D1N3-2]%[ma:5d3ebb90]
% nguon: Toán tư duy lớp 4 – Bài 3 SGK Toán 4 KNTT
	Không cần tính, cho biết kết quả mỗi phép tính sau là số chẵn hay số lẻ:\\
	a) $135 + 247$;\qquad b) $482 - 159$;\qquad c) $25 \times 46$.
	\loigiai{
		Nhớ: lẻ $+$ lẻ $=$ chẵn; chẵn $-$ lẻ $=$ lẻ; tích có một thừa số chẵn thì là số chẵn.\\
		a) $135$ lẻ, $247$ lẻ nên tổng là số chẵn.\\
		b) $482$ chẵn, $159$ lẻ nên hiệu là số lẻ.\\
		c) $46$ chẵn nên tích là số chẵn.
	}
\end{bt}

% nguồn: 2026-11-18_lop4_chu-de-1-tu-luan (nhập 18/11/2026)
\begin{bt}%[4D1H3-2]%[ma:947e7e97]
% nguon: Toán tư duy lớp 4 – Bài 3 SGK Toán 4 KNTT
	Tổng của $5$ số lẻ bất kì là số chẵn hay số lẻ? Vì sao?
	\loigiai{
		Ghép $4$ số lẻ đầu thành $2$ cặp, mỗi cặp lẻ $+$ lẻ $=$ chẵn, nên tổng $4$ số đó là số chẵn.\\
		Cộng thêm số lẻ thứ năm: chẵn $+$ lẻ $=$ lẻ.\\
		Vậy tổng của $5$ số lẻ là số lẻ.
	}
\end{bt}

% nguồn: 2026-11-18_lop4_chu-de-1-tu-luan (nhập 18/11/2026)
\begin{bt}%[4D1V3-2]%[ma:a76e20d1]
% nguon: Toán tư duy lớp 4 – Bài 3 SGK Toán 4 KNTT
	Bạn Bình tính tổng $1 + 2 + 3 + \dots + 15$ được kết quả $121$. Không cần tính lại tổng, hãy giải thích vì sao Bình tính sai.
	\loigiai{
		Từ $1$ đến $15$ có các số lẻ $1, 3, 5, 7, 9, 11, 13, 15$, tức là $8$ số lẻ.\\
		Tổng của $8$ số lẻ (ghép thành $4$ cặp lẻ $+$ lẻ) là số chẵn; cộng thêm các số chẵn vẫn được số chẵn.\\
		Vậy tổng phải là số chẵn, mà $121$ là số lẻ nên Bình tính sai. (Tổng đúng là $120$.)
	}
\end{bt}

% nguồn: 2026-11-18_lop4_chu-de-1-tu-luan-bo-sung (nhập 18/11/2026)
\begin{bt}%[4D1H3-2]%[ma:4178dc0e]
% nguon: Toán tư duy lớp 4 – Bài 3 SGK Toán 4 KNTT
	Hiệu của hai số tự nhiên là $37$. Hai số đó có thể cùng là số chẵn không? Có thể cùng là số lẻ không? Vì sao?
	\loigiai{
		Chẵn $-$ chẵn $=$ chẵn; lẻ $-$ lẻ $=$ chẵn.\\
		Hiệu là $37$, một số lẻ, nên hai số không thể cùng chẵn, cũng không thể cùng lẻ.\\
		Vậy hai số đó phải có một số chẵn và một số lẻ.
	}
\end{bt}
